\documentclass{article}
\usepackage{amsmath,amssymb,amsthm}
\usepackage{gb4e}

\newtheorem{definition}{Definition}
\newtheorem{proposition}{Proposition}
\newtheorem{prediction}{Prediction}
\newcommand{\sem}[1]{[\![#1]\!]}
\newcommand{\lam}[1]{\lambda #1 \,.\,}

\title{Determiners as Relations: A Toy Fragment}
\author{Anonymous}

\begin{document}
\maketitle

\section{Introduction}

This note develops a small extensional fragment of English containing
proper names, intransitive verbs, common nouns and the determiners
\textit{every}, \textit{some} and \textit{no}, and checks some of its
predictions about sentences like those in (\ref{ex:ann})--(\ref{ex:some-phon}).

\begin{exe}
  \ex\label{ex:ann} Ann sneezed.
  \ex\label{ex:every} Every linguist sneezed.
  \ex\label{ex:some} Some linguist sneezed.
  \ex\label{ex:no} No linguist sneezed.
  \ex\label{ex:every-phon} Every phonologist sneezed.
  \ex\label{ex:some-phon} Some phonologist sneezed.
\end{exe}

\section{The fragment}

We work with an extensional type theory with basic types $e$ and $t$, where
$D_e$ is the domain of entities and $D_t = \{0, 1\}$. For types $\sigma,
\tau$, $D_{\langle\sigma,\tau\rangle}$ is the set of functions from
$D_\sigma$ to $D_\tau$. A \emph{model} fixes $D_e$ together with the
interpretation of the non-logical vocabulary (\textit{Ann}, \textit{linguist},
\textit{phonologist}, \textit{sneezed}). We freely identify a function $f \in
D_{\langle e,t\rangle}$ with the set $\{x \in D_e \mid f(x) = 1\}$.

\paragraph{Syntax.} We assume the following binary-branching LFs, where
(\ref{ex:ann}) is $[_S\ \text{Ann}\ \text{sneezed}]$, and a quantified sentence
such as (\ref{ex:every}) is $[_S\ [_{DP}\ \text{every}\ \text{linguist}]\
\text{sneezed}]$; (\ref{ex:some})--(\ref{ex:some-phon}) are parallel.

\begin{definition}[Lexicon]\label{def:lexicon}
\begin{align*}
  \sem{\textit{Ann}} &= \mathbf{ann} \in D_e\\
  \sem{\textit{linguist}} &= \lam{x \in D_e} x \text{ is a linguist}\\
  \sem{\textit{phonologist}} &= \lam{x \in D_e} x \text{ is a phonologist}\\
  \sem{\textit{sneezed}} &= \lam{x \in D_e} x \text{ sneezed}\\
  \sem{\textit{every}} &= \lam{f \in D_{\langle e,t\rangle}} \lam{g \in D_{\langle e,t\rangle}} \forall x \in D_e\,[f(x) = 1 \rightarrow g(x) = 1]\\
  \sem{\textit{some}} &= \lam{f \in D_{\langle e,t\rangle}} \lam{g \in D_{\langle e,t\rangle}} \exists x \in D_e\,[f(x) = 1 \wedge g(x) = 1]\\
  \sem{\textit{no}} &= \lam{f \in D_{\langle e,t\rangle}} \lam{g \in D_{\langle e,t\rangle}} \neg\exists x \in D_e\,[f(x) = 1 \wedge g(x) = 1]
\end{align*}
\end{definition}

\begin{definition}[Composition rules]\label{def:rules}
\textbf{Terminal Nodes (TN):} if $\alpha$ is a terminal node, $\sem{\alpha}$
is specified in the lexicon. \textbf{Functional Application (FA):} if
$\alpha$ is a branching node with daughters $\beta$ and $\gamma$, and
$\sem{\beta} \in D_{\langle\sigma,\tau\rangle}$ and $\sem{\gamma} \in
D_\sigma$, then $\sem{\alpha} = \sem{\beta}(\sem{\gamma})$.
\end{definition}

We restrict attention to models that respect the obvious lexical relation
between the two nouns:

\begin{definition}[Meaning postulate]\label{def:mp}
For all $x \in D_e$: if $\sem{\textit{phonologist}}(x) = 1$ then
$\sem{\textit{linguist}}(x) = 1$.
\end{definition}

\begin{definition}[Entailment]\label{def:entail}
A sentence $S$ \emph{entails} a sentence $S'$ iff $\sem{S'} = 1$ in every
model (respecting Definition~\ref{def:mp}) in which $\sem{S} = 1$. Two
sentences are \emph{contradictory} iff there is no such model in which both
are true.
\end{definition}

\section{Truth conditions}

For (\ref{ex:ann}), FA applies once:
$\sem{(\ref{ex:ann})} = \sem{\textit{sneezed}}(\sem{\textit{Ann}}) = 1$ iff
$\mathbf{ann}$ sneezed.
The derivation for (\ref{ex:every}) is as follows.
\begin{align*}
  \sem{(\ref{ex:every})}
  &= \sem{[_{DP}\ \text{every linguist}]}(\sem{\textit{sneezed}}) && \text{by FA}\\
  &= \sem{\textit{every}}(\sem{\textit{linguist}})(\sem{\textit{sneezed}}) && \text{by FA}\\
  &= 1 \text{ iff } \forall x \in D_e\,[\sem{\textit{linguist}}(x) = 1 \rightarrow \sem{\textit{sneezed}}(x) = 1] && \text{by TN, $\beta$-reduction}\\
  &= 1 \text{ iff every linguist sneezed}
\end{align*}
Entirely analogous derivations give $\sem{(\ref{ex:some})} = 1$ iff $\exists
x \in D_e\,[x \text{ is a linguist} \wedge x \text{ sneezed}]$, and
$\sem{(\ref{ex:no})} = 1$ iff $\exists x \in D_e\,[x \text{ is a linguist}
\wedge \neg(x \text{ sneezed})]$.

\section{Determiner properties}

\begin{definition}[Conservativity]\label{def:cons}
A determiner denotation $\delta$ is \emph{conservative} iff for all $A, B$:
$\delta(A)(B) = \delta(A)(A \cap B)$.
\end{definition}

\begin{definition}[Monotonicity]\label{def:mon}
A determiner denotation $\delta$ is \emph{left downward monotone}
(${\downarrow}\mathrm{MON}$) iff whenever $\delta(A)(B) = 1$ and
$A' \subseteq A$, then $\delta(A')(B) = 1$. It is \emph{right upward monotone}
($\mathrm{MON}{\uparrow}$) iff whenever $\delta(A)(B) = 1$ and
$B \subseteq B'$, then $\delta(A)(B') = 1$. Right downward monotonicity
($\mathrm{MON}{\downarrow}$) is defined analogously.
\end{definition}

\begin{proposition}\label{prop:cons}
$\sem{\textit{every}}$, $\sem{\textit{some}}$ and $\sem{\textit{no}}$ are
all conservative.
\end{proposition}

\begin{proposition}\label{prop:every-mon}
$\sem{\textit{every}}$ is ${\downarrow}\mathrm{MON}$ and $\mathrm{MON}{\uparrow}$.
\end{proposition}

\begin{proposition}\label{prop:no-mon}
$\sem{\textit{no}}$ is ${\downarrow}\mathrm{MON}$ and $\mathrm{MON}{\downarrow}$.
\end{proposition}

\begin{proposition}\label{prop:some-mon}
$\sem{\textit{some}}$ is ${\downarrow}\mathrm{MON}$.
\end{proposition}

\section{Predictions}

The determiner properties above translate directly into predictions about
the sentences in (\ref{ex:ann})--(\ref{ex:some-phon}).

\begin{prediction}\label{pred:every-phon}
(\ref{ex:every}) entails (\ref{ex:every-phon}).
\end{prediction}

This follows from Proposition~\ref{prop:every-mon} together with
Definition~\ref{def:mp}. By the same reasoning, Proposition~\ref{prop:some-mon}
predicts that (\ref{ex:some}) entails (\ref{ex:some-phon}).

\begin{prediction}\label{pred:some-phon-up}
(\ref{ex:some-phon}) entails (\ref{ex:some}).
\end{prediction}

\begin{prediction}\label{pred:instantiation}
In any model in which $\mathbf{ann}$ is a linguist, (\ref{ex:every}) entails
(\ref{ex:ann}).
\end{prediction}

\begin{prediction}\label{pred:contradictory}
(\ref{ex:some}) and (\ref{ex:no}) are contradictory.
\end{prediction}

Finally, the fragment captures the intuition that universal statements
carry existential commitment: (\ref{ex:every}) entails (\ref{ex:some}), since
if every linguist sneezed then, in particular, some linguist did.

\end{document}
